\section{Learning path: six steps from one qubit to multi-stage circuits}
\hypertarget{learning-path}{}

Each step below uses only ideas from the steps before it. You do not need to
finish the theory chapters before starting: build the starter circuit first,
predict what it should do, then read the matching explanation.

\begin{center}
\begin{tikzpicture}[font=\scriptsize,
  rung/.style={draw=QpuBlue,fill=QpuBlue!8,rounded corners=3pt,text width=2.3cm,
    align=center,minimum height=1.15cm,inner sep=3pt},
  >={Latex[length=2mm]}]
  \node[rung] (s1) at (0,0) {\textbf{1. Single}\\\textbf{qubit}\\X, H, Z, M};
  \node[rung] (s2) at (2.95,0) {\textbf{2. Controlled}\\\textbf{gates}\\CNOT, CZ, SWAP};
  \node[rung] (s3) at (5.9,0) {\textbf{3. Reversible}\\\textbf{logic}\\CCNOT, AND, XOR};
  \node[rung] (s4) at (8.85,0) {\textbf{4. Entanglement}\\\strut\\Bell, kickback};
  \node[rung] (s5) at (11.8,0) {\textbf{5. Arithmetic}\\\strut\\adders, parity};
  \node[rung] (s6) at (14.75,0) {\textbf{6. Multi-stage}\\\textbf{circuits}\\children, oracles};
  \foreach \a/\b in {s1/s2,s2/s3,s3/s4,s4/s5,s5/s6} \draw[->,thick,QpuNavy] (\a)--(\b);
\end{tikzpicture}
\end{center}

\begin{beginnerbox}[How to use each step]
\begin{enumerate}
  \item Open \textbf{Circuit builder} and choose the starter circuit named in
  the step from \textbf{Load a starter circuit}. Each card shows its step
  number.
  \item Before running it, write down what you expect to measure.
  \item Choose \textbf{Step gate} repeatedly and watch the highlighted column
  on the canvas and the \textbf{Execution log}.
  \item Compare your prediction with the result, then read the listed guide
  section.
  \item Move on when you can answer the step's checkpoint question without
  looking.
\end{enumerate}
\end{beginnerbox}

\subsection{Step 1: one qubit at a time}

\textbf{Learn:} a qubit starts in $\ket0$; X flips it; H creates an equal
superposition; Z changes phase without changing immediate probabilities;
MEASURE (the meter on the qubit wire) samples a classical 0 or 1 and marks that time on the \textbf{c} lane. The qubit can still be used afterward.

\textbf{Starter circuits:} \emph{Superposition coin}, then
\emph{Phase interference (H-Z-H)}.

\textbf{Read:} \emph{First circuit} and \emph{Randomness versus interference}
in Getting Started; \emph{Single-qubit primitive gates} in the Theory Guide.

\textbf{Checkpoint:} why does H followed by M give 0 or 1 at random, while
H, Z, H followed by M always gives 1?

\subsection{Step 2: controlled gates}

\textbf{Learn:} a controlled gate reads one wire (the \emph{control}) and
changes another (the \emph{target}). The control is never modified by CNOT.
In the builder, the control comes from \textbf{Control A} and the target from
\textbf{Target particle}.

\textbf{Starter circuit:} \emph{CNOT demo}. Change \textbf{q0 start} between
\code{0p} and \code{1p} and rerun.

\textbf{Read:} \emph{Controlled and two-qubit primitive gates} in the Theory
Guide.

\textbf{Checkpoint:} in \emph{CNOT demo}, which wire can change, and under what
condition?

\subsection{Step 3: reversible Boolean logic}

\textbf{Learn:} quantum gates cannot throw information away, so reversible
logic keeps its inputs and writes the answer into a separate
\emph{target} (also called workspace). Toffoli (CCNOT) and the derived AND gate
both compute $t'=t\oplus(A\land B)$: A and B are read and left unchanged, and
$t$ is the output wire. Start $t$ at 0 and it ends holding $A\land B$.

\textbf{Starter circuits:} \emph{CCNOT demo}, then \emph{Half adder}.

\textbf{Read:} \emph{Classical computation foundations} and \emph{Worked
example: reversible half-adder} in Getting Started; \emph{Derived logical
gates} in the Theory Guide.

\textbf{Checkpoint:} if the target of AND starts at 1 instead of 0, what does
it hold afterward?

\subsection{Step 4: entanglement and phase}

\textbf{Learn:} a controlled gate whose control is in superposition can
correlate two wires so that neither has a separate state (a Bell state). Phase
can also flow \emph{backward} from target to control (phase kickback).

\textbf{Starter circuits:} \emph{Bell state}, then \emph{Phase kickback}.

\textbf{Read:} \emph{Multiple qubits} (Entanglement) and \emph{Measurement,
probability, and truth tables} in the Theory Guide; \emph{Phase kickback} in
Examples and Troubleshooting.

\textbf{Checkpoint:} after \emph{Bell state}, if q0 measures 1, what will q1
measure? Why is that not the same as copying q0?

\subsection{Step 5: arithmetic}

\textbf{Learn:} addition is built from the step 3 pieces. Sum bits are parity
(XOR), carry bits are majority (a combination of ANDs). A comparator or parity
checker uses the same pieces with a different output.

\textbf{Starter circuits:} \emph{Parity check}; then compile the bundled
\code{SingleBitFullAdder} and \code{TwoBitFullAdder} buttons in
\textbf{Compile text into a circuit}.

\textbf{Read:} \emph{Advanced worked circuits} in Examples and Troubleshooting
(full adder, two-bit adder, comparators, parity).

\textbf{Checkpoint:} why does a full adder need three AND-style gates for its
carry, but only XOR-style gates for its sum?

\subsection{Step 6: multi-stage circuits}

\textbf{Learn:} larger circuits are built in stages: compute a helper value,
use it, then \emph{uncompute} it so workspace returns to 0. Child processes
(\code{RUNCHILD}) and custom gates reuse whole circuits as single steps.
Oracle algorithms combine all of the above.

\textbf{Starter circuit:} \emph{Grover search (2 qubits)}; then the
multi-stage listings in \emph{Advanced worked circuits}.

\textbf{Read:} \emph{Process composition and return values} in the Language
Reference; \emph{Custom gates} in Examples and Troubleshooting.

\textbf{Checkpoint:} after an uncompute stage, what value should every
temporary wire hold, and how can a truth-table test prove it?
